\documentclass[
  draft,
  linenumbers
]{agujournal2025}

\usepackage{booktabs}
\usepackage{amsmath}
\usepackage{siunitx}

\graphicspath{{figures/}}

\newcommand{\coTwo}{CO$_2$}
\newcommand{\todo}[1]{\textbf{[TODO: #1]}}

\begin{document}

\journalname{Water Resources Research}

\title{End-to-End Uncertainty Quantification of Geologic \coTwo{} Storage in Faulted Reservoirs}

\authors{Shaowen Mao\affil{1}, Llu\'{\i}s Sal\'o-Salgado\affil{2},
Hannah Lu\affil{3}, Lisa Lun\affil{4}, Christie M. Rogers\affil{4},
Rodrick D. Myers\affil{4}, Nathan W. Harkins\affil{4},
Patricia Ascanio-Pellon\affil{4}, and Ruben Juanes\affil{1,5}}

\affiliation{1}{Department of Civil and Environmental Engineering,
Massachusetts Institute of Technology, Cambridge, MA, USA}
\affiliation{2}{Department of Earth, Environmental, and Planetary Sciences,
University of Tennessee, Knoxville, TN, USA}
\affiliation{3}{The Oden Institute for Computational Engineering and Sciences,
The University of Texas at Austin, Austin, TX, USA}
\affiliation{4}{ExxonMobil Technology and Engineering Company, Spring, TX, USA}
\affiliation{5}{Department of Earth, Atmospheric and Planetary Sciences,
Massachusetts Institute of Technology, Cambridge, MA, USA}

\authoraddr{Corresponding author: Shaowen Mao (\todo{email})}

\authorrunninghead{Mao et al.}
\titlerunninghead{End-to-end UQ of \coTwo{} storage in faulted reservoirs}

\keypoints
  {An end-to-end workflow propagates fault uncertainty from geologic scenarios to field-scale \coTwo{} migration}
  {Hierarchical sampling preserves joint fault properties while reducing 162 geologies to 1,620 simulations}
  {The ensemble quantifies containment, fault and top-seal entry, vertical migration, and \coTwo{} phase partitioning}

\maketitle

\begin{abstract}
Faults are ubiquitous in the subsurface and can act as seals or flow conduits. In geologic carbon sequestration (GCS), uncertainty in how injected \coTwo{} interacts with faults hinders decisions about migration and containment. We develop an end-to-end uncertainty quantification workflow for \coTwo{} migration along and across faults in realistic geologic settings. First, the stochastic fault-property model PREDICT generates conditional distributions of fault permeability from local stratigraphy and fault throw. Second, a hierarchical sampling algorithm constructs representative realizations of upscaled three-dimensional fault properties. Third, high-performance multiphase reservoir simulations model megatonne-scale injection and migration over 1,000 years. We apply the workflow to an offshore Texas field represented by a roughly two-million-cell geologic model containing a three-dimensional growth fault that partially offsets the top seal. We construct 162 geologic scenarios by varying key structural and stratigraphic parameters. For each scenario, PREDICT generates fine-scale fault-core architectures and conditional distributions of upscaled fault permeability within six throw windows. Hierarchical sampling constructs 10 representative three-dimensional fault-property realizations per geology, yielding 1,620 field-scale simulations with spatially variable permeability, porosity, capillary pressure, and relative permeability. The ensemble quantifies uncertainty in the fractions of injected \coTwo{} retained within the storage reservoir and entering the fault and top seal, maximum upward migration, and dissolved-versus-free \coTwo{} partitioning. The framework propagates uncertainty from geologic characterization and stochastic fault-property prediction to field-scale migration and containment, supporting GCS site selection, injection management, and monitoring design.

\begin{plainlanguagesummary}
Faults can either impede or promote the migration of injected \coTwo{}, but their properties cannot be observed everywhere or represented by one deterministic value. We combine stochastic fault-property prediction, uncertainty-aware sampling, and field-scale reservoir simulation to carry this uncertainty from geologic interpretation through 1,000 years of modeled storage. Applied to an offshore Texas field, the workflow converts 162 plausible geologies into 1,620 reproducible simulations. The ensemble is designed to show how uncertain fault properties affect storage-reservoir retention, entry into the fault and top seal, upward migration, and the division of \coTwo{} between dissolved and free phases. These results can inform where to store \coTwo{}, how to operate injection, and what monitoring should target.
\end{plainlanguagesummary}
\end{abstract}

% !TEX root = ../manuscript.tex

\section{Introduction}

Faults are ubiquitous in sedimentary basins and may either impede or promote
subsurface fluid migration. In geologic carbon sequestration, this dual role
creates uncertainty in whether injected \coTwo{} remains in the intended
storage interval, enters a fault, migrates toward or into the top seal, or
moves upward along the fault. Quantifying these possibilities is important for
site selection, injection management, and monitoring design.

The effective flow behavior of a fault emerges from unresolved internal
structure, including juxtaposed sand and clay units, stochastic clay-smear
continuity, burial and faulting history, and directional permeability.
Representing the fault by one deterministic property field can therefore
understate both the range and the structured, multimodal character of
uncertainty relevant to reservoir-scale forecasts. PREDICT provides a
physics-informed stochastic representation of fault-core architecture and
effective permeability for normally consolidated sediments
\cite{SaloSalgado2025}. Previous work has also demonstrated the importance of
propagating fault-property uncertainty into predictions of \coTwo{} migration
and fault response \cite{Lu2025}.

The remaining challenge is computational scale. An end-to-end assessment must
span plausible geologies, stochastic fault realizations within each throw
window, uncertain relationships among windows, spatial variability along
strike, upscaled multiphase properties, and long-term field-scale simulation.
Direct enumeration of this combined space is infeasible, whereas aggressive
averaging can remove permeability-tensor dependence, distinct probability
modes, and capillary barriers that control migration.

We develop an end-to-end uncertainty quantification workflow with three
computational stages: stochastic fault-property prediction, hierarchical
sampling of three-dimensional fault properties, and high-performance
multiphase reservoir simulation. The sampling hierarchy separates geologic
uncertainty, within-window PREDICT stochasticity, and cross-window coupling
uncertainty. PREDICT produces window-specific distributions but does not
identify a true cross-window joint law; consequently, distributional
similarity is used to construct transparent candidate coupling assumptions,
not to claim measured cross-window correlation.

We apply the workflow to a megatonne-scale offshore Texas storage case with a
three-dimensional growth fault that partially offsets the top seal. The study
uses 162 geologic scenarios and 10 sampled fault-property cases per geology,
yielding 1,620 reservoir simulations over a 1,000-year horizon. We quantify
storage-reservoir retention, entry into the fault and top seal, maximum upward
migration, and partitioning between dissolved and free \coTwo{}. The objectives
are to determine how uncertainty propagates from geologic characterization to
these containment outcomes, identify the dominant controls, and evaluate the
implications for storage decisions. \todo{Add a concise statement of the main
quantitative findings after the production ensemble is complete.}

% !TEX root = ../manuscript.tex

\section{Offshore Texas Field Case and Geologic Uncertainty}

\subsection{Field Setting and Geologic Model}

We apply the workflow to an offshore Texas storage setting represented by a
three-dimensional geologic model with approximately two million grid cells.
The model contains a growth fault that offsets the storage interval and
partially offsets the top seal. This configuration permits migration both
along the fault and across juxtaposed units, making the spatially variable
fault properties central to predictions of containment.

\todo{Describe the storage interval, top-seal stratigraphy, model dimensions,
grid resolution, fault geometry, boundary conditions inherited from the field
model, and the evidence supporting the selected geologic interpretation. Add a
field-location and model-geometry figure.}

\subsection{Fault Discretization and Throw Windows}

The reservoir-scale fault representation contains 87 cells along strike and
six throw windows, denoted W1--W6. Each throw window has a different local
stratigraphic juxtaposition and therefore a distinct conditional distribution
of fault properties. The six window properties are assigned independently at
each along-strike location unless a sampled fault-wide or grouped case imposes
a shared low/high state. This $6\times87$ representation connects the
window-scale PREDICT libraries to the three-dimensional fault surface used by
the reservoir model.

\todo{Define the throw ranges associated with W1--W6 and show how the window
and along-strike indices map onto the reservoir grid.}

\subsection{Geologic Scenario Design}

Geologic uncertainty is the outer level of the hierarchy. We consider six
layer-thickness structures spanning low, medium, and high sand proportions,
each represented with uniform and nonuniform layer thicknesses. Each structure
is combined factorially with three faulting depths, three clay-volume
fractions in sand, and three clay-volume fractions in clay, producing

\begin{equation}
6\times3\times3\times3 = 162
\end{equation}

geologic scenarios. All scenarios receive equal treatment in the downstream
workflow; no clustering or representative-scenario reduction is applied at
this outer level.

\todo{Add the final parameter values, units, geologic rationale, and a table
that maps scenario identifiers to the complete factorial design.}

% !TEX root = ../manuscript.tex

\section{End-to-End Uncertainty Quantification Workflow}

\subsection{Workflow Overview}

The computational pipeline has three stages: stochastic fault-property
prediction with PREDICT, hierarchical sampling of three-dimensional fault
properties, and field-scale multiphase reservoir simulation. Within the
sampling stage, uncertainty is organized into three levels: (1) geologic
scenario uncertainty, (2) within-window PREDICT stochasticity, and (3)
cross-window coupling uncertainty. Distinguishing the pipeline stages from the
uncertainty levels clarifies which quantities are predicted, which are sampled,
and which are propagated through the reservoir model.

\todo{Add an end-to-end workflow figure that tracks inputs, empirical
libraries, sampled $6\times87$ property fields, upscaled multiphase
properties, reservoir simulations, and quantities of interest.}

\subsection{Stochastic Fault-Property Prediction With PREDICT}

For each geology and throw window, PREDICT generates stochastic fine-scale
fault-core architectures conditioned on the local stratigraphy, fault throw,
and geologic parameters. Each accepted realization yields an upscaled joint
permeability vector $(k_{xx},k_{yy},k_{zz})$. We generate 2,000 accepted
realizations for every geology-window combination. A separate convergence
analysis established that 2,000 realizations bring the component-distribution
distances close to the empirical floor between independent 20,000-realization
reference ensembles.

Adjacent layers of the same lithology are collapsed before fault-core
generation. Clay-smear placement uses the cell-union probability-of-smear
rule: each clay source produces an active cell mask, a cell remains clay if
any source occupies it, and overlapping sources are assigned a parent source
after clay presence is established. This ordering prevents competition among
overlapping sources from converting a potentially clay-occupied cell back to
sand.

\todo{Add the governing PREDICT relations, the material-placement algorithm,
the effective-property upscaling equations, and the final validation evidence
for the cell-union probability-of-smear option.}

\subsection{Hierarchical Sampling of Fault Properties}

\subsubsection{Within-window joint permeability states}

Each PREDICT realization is represented by
\begin{equation}
\mathbf{y}^{(m)} = \left(\log_{10} k_{xx}^{(m)},
\log_{10} k_{yy}^{(m)},\log_{10} k_{zz}^{(m)}\right).
\end{equation}
Joint permeability clusters are detected by $k$-medoids in physical
log-permeability space using Euclidean distance with one log unit as the
component scale. Candidate cluster counts are $K=1,\ldots,5$. Multi-cluster
solutions must satisfy a minimum cluster fraction of 5\%, and the valid
solution with the largest mean silhouette is selected. If its silhouette is
below 0.25, the window is treated as single-cluster.

For interpretation, local percentile ranks $r_{xx}$, $r_{yy}$, and $r_{zz}$
are computed within each window. The joint rank score is
\begin{equation}
s^{(m)} = \frac{r_{xx}^{(m)}+r_{yy}^{(m)}+r_{zz}^{(m)}}{3}.
\end{equation}
Clusters are ordered from low to high by their median joint rank score. Low
and high state libraries each contain 20\% of the full window library.
Complete extreme clusters are included first; if the target size would be
exceeded, only the required lowest- or highest-score samples from the boundary
cluster are retained. This cluster-aware quantile construction preserves
multimodal structure while maintaining a fixed state mass.

Each state medoid is an actual PREDICT realization. The local perturbation
pool contains samples nearest to the medoid within the intersection of the
state library and the medoid-containing cluster. Its size is the smaller of
the available candidate count and the larger of 25 samples or 10\% of that
candidate set. The state-wide pool is the complete low or high state library.
Independent samples are drawn from the complete window library. In every
case, $k_{xx}$, $k_{yy}$, and $k_{zz}$ are selected together from one actual
realization.

\subsubsection{Cross-window similarity and coupling cases}

For each geology, pairwise differences between the six joint permeability
distributions are quantified with normalized multivariate energy distance.
The normalization compares between-window distributional separation with the
characteristic within-window spread. Windows form a similar pair when the
normalized distance is no greater than 0.25. Optional bootstrap resampling
tests the stability of this decision but is not required for the production
workflow because the empirical window libraries were independently
convergence tested.

All 203 partitions of six windows are enumerated. A partition is admissible
only when every pair within every non-singleton group satisfies the similarity
threshold, preventing chain-link grouping. Among admissible partitions, the
workflow first selects the smallest group count and then breaks ties using the
smallest aggregate within-group distance. Singleton windows remain independent
because the distributions do not support coupling them to another window.

Ten multiple-window permeability cases are constructed for each geology: two
independent cases, four fault-wide low/high cases using local or state-wide
pools, and four grouped low/high cases. Non-singleton groups share low/high
state assignments, whereas singleton windows remain independently sampled.
When more than two non-singleton groups are present, candidate binary group
splits are scored by the difference between mean opposite-side and same-side
group distance; the highest-scoring split and its complement are retained.

\subsection{Along-Strike Sampling, Exact Replay, and Multiphase Upscaling}

For each of the 10 cases, the prescribed six-window design is sampled
independently at all 87 along-strike locations. Deterministic seeds and the
accepted PREDICT realization identifiers make every selected permeability
vector exactly replayable. Replay reconstructs the corresponding fine-scale
fault-core maps needed to upscale porosity, capillary pressure, and relative
permeability without storing every fine-scale realization during the original
PREDICT campaign.

The rigorous capillary-pressure workflow uses invasion percolation on each
replayed fault-core realization and retains spatially variable saturation
endpoints. The reservoir-input workflow can either preserve every slice-level
capillary-pressure curve or reduce curves within physically distinct entry-
pressure branches using branch-specific medoids. Dynamic relative-permeability
upscaling is performed for a representative saturation-endpoint realization
in each window and mapped consistently to the retained saturation regions.

\todo{Complete the equations and numerical details for porosity,
invasion-percolation capillary pressure, effective irreducible water
saturation, dynamic relative permeability, branch detection, medoid
selection, endpoint normalization, and quality-control criteria.}

% !TEX root = ../manuscript.tex

\section{Field-Scale Multiphase Simulation and Analysis}

\subsection{Reservoir Model and Injection Scenario}

High-performance multiphase simulations model megatonne-scale \coTwo{}
injection and subsequent migration over 1,000 years in the approximately
two-million-cell offshore Texas model. The model resolves the storage
reservoir, top seal, and three-dimensional growth fault so that migration can
be tracked within the reservoir and both along and across the fault.

\todo{Report the simulator and nonlinear-solver configuration, fluid model,
initial and boundary conditions, injection rate and duration, well controls,
rock and fluid properties outside the fault, grid statistics, timestep
strategy, convergence tolerances, and computing platform.}

\subsection{Fault-Property Mapping and Ensemble Design}

Each of the 162 geologic scenarios contributes 10 sampled three-dimensional
fault-property cases, giving
\begin{equation}
162\times10 = 1{,}620
\end{equation}
field-scale simulations. Every case contains spatially variable fault
permeability and porosity over the $6\times87$ fault representation, together
with saturation-region assignments for the selected capillary-pressure and
relative-permeability curves. The permeability case design is deterministic
for a given geology and case type, while draws within the prescribed sampling
pools vary reproducibly along strike.

\todo{Define the final reservoir-ready file schema, property-to-grid mapping,
saturation-region treatment, random-seed policy, restart protocol, and
simulation acceptance criteria.}

\subsection{Migration and Containment Quantities of Interest}

The simulation ensemble is evaluated using quantities of interest that link
fault behavior to storage security:
\begin{enumerate}
  \item the fraction of injected \coTwo{} retained within the intended storage
  reservoir;
  \item the fraction entering the fault;
  \item the fraction entering the top seal;
  \item the maximum upward migration distance; and
  \item the partitioning of \coTwo{} between dissolved and free phases.
\end{enumerate}

\todo{Provide exact spatial masks, denominators, evaluation times, phase and
component definitions, and numerical thresholds for every quantity of
interest. Distinguish instantaneous maxima from values evaluated at the end of
injection and after 1,000 years.}

\subsection{Ensemble Analysis}

The ensemble analysis compares outcomes across geologic scenarios, within-
window stochastic states, and cross-window coupling assumptions. This
hierarchical labeling enables both total predictive ranges and conditional
comparisons that identify which uncertainty source controls each migration or
containment metric.

\todo{Specify the statistical estimands, uncertainty-decomposition method,
matched comparisons among the 10 case types, treatment of failed simulations,
and confidence or sensitivity diagnostics.}

% !TEX root = ../manuscript.tex

\section{Results}

\subsection{Geologic Controls on PREDICT Fault Properties}

\todo{Report matched-effect estimates for sand proportion, thickness
uniformity, faulting depth, clay volume fraction in sand, and clay volume
fraction in clay. Separate directional effects on $k_{xx}$, $k_{yy}$, and
$k_{zz}$ and connect them to the fault-material and tensor-rotation mechanisms
in PREDICT.}

\subsection{Hierarchical Reduction and Three-Dimensional Fault Cases}

\todo{Summarize cluster-count frequencies, state-library composition, window
similarity-group structures, pairwise co-grouping frequencies, and their
dependence on geology. Show how independent, fault-wide low/high, and grouped
low/high designs translate into $6\times87$ fault-property fields. Emphasize
that the reduction retains actual joint realizations and important multimodal
behavior.}

\subsection{Upscaled Multiphase Fault Properties}

\todo{Present spatial variation in porosity, capillary-pressure entry levels,
effective irreducible water saturation, and relative permeability. Compare the
full slice-level and entry-pressure-branch representations and document the
final choice used in the reservoir ensemble.}

\subsection{Field-Scale \coTwo{} Migration and Containment}

\todo{Report storage-reservoir retention, entry into the fault and top seal,
maximum upward migration, and dissolved-versus-free phase partitioning over
the injection and post-injection periods. Lead with the principal physical and
decision-relevant findings rather than intermediate workflow diagnostics.}

\subsection{Contributions of the Uncertainty Hierarchy}

\todo{Quantify how much outcome variability is associated with geology,
within-window PREDICT stochasticity, and cross-window coupling assumptions.
Identify interactions and the conditions responsible for the most secure and
least secure migration outcomes.}

% !TEX root = ../manuscript.tex

\section{Discussion}

\subsection{Implications for Site Selection, Operations, and Monitoring}

\todo{Translate the ensemble results into implications for screening faulted
storage sites, selecting injection rates and locations, and prioritizing
monitoring of the fault, top seal, and upward-migration pathways. Distinguish
robust implications from field-specific observations.}

\subsection{Value and Interpretation of Hierarchical Reduction}

The full uncertainty space combines 162 geologies, six window-specific
empirical libraries, multiple cross-window assignments, and along-strike
variability. Exhaustive reservoir simulation is therefore infeasible. The
hierarchical reduction retains actual joint permeability realizations,
dominant multimodal structure, and explicit coupling assumptions rather than
replacing the empirical distributions with component-wise averages or one
parametric model.

Window similarity groups identify distributions that can reasonably share a
state assignment in candidate coupled cases. They do not establish that
permeability realizations are physically correlated across windows. This
distinction is necessary because PREDICT is run independently by window and
does not identify the true cross-window joint distribution.

\subsection{Representation of Multiphase Fault Properties}

\todo{Discuss why spatially variable capillary pressure cannot generally be
replaced by a simple average of sand and clay reference curves, when entry-
pressure-branch medoids provide an acceptable simplification, and why a
representative relative-permeability curve is more defensible after endpoint-
consistent scaling. Relate these choices to reservoir-solver robustness.}

\subsection{Limitations and Transferability}

\todo{Address PREDICT structural assumptions, validation of clay-smear
placement, state and similarity thresholds, absence of observed cross-window
dependence, finite sampling and reservoir-simulation budgets, model and fluid-
property uncertainty outside the fault, and transferability beyond the
offshore Texas field.}

% !TEX root = ../manuscript.tex

\section{Conclusions}

We developed an end-to-end workflow that propagates uncertainty from geologic
scenario design and stochastic fault-core prediction through hierarchical
three-dimensional property sampling and field-scale multiphase simulation.
For 162 geologic scenarios, the workflow constructs 10 reproducible fault-
property cases per geology while preserving actual joint PREDICT permeability
realizations, along-strike variability, and explicit cross-window coupling
assumptions. Exact replay connects selected effective permeability tensors to
the fine-scale properties needed for porosity, capillary-pressure, and
relative-permeability upscaling.

\todo{Replace this provisional conclusion with three to five quantitative
findings: the dominant geologic control, the contribution of stochastic and
coupling uncertainty, the range of migration and containment outcomes, the
most decision-relevant security metric, and the practical implication for GCS
site selection, injection management, or monitoring design.}


\section*{Acknowledgments}
\todo{Add funding, computing resources, technical contributions, and author contributions.}

\section*{Conflict of Interest}
\todo{Add the final conflict-of-interest disclosure for all authors.}

\section*{Open Research Statement}
The PREDICT source code and the reproducible analysis workflow are maintained at \url{https://github.com/ShaowenMao/predict_shaowen}. \todo{Deposit the manuscript version of the software and derived data in archival repositories, add DOIs, licenses, and citations, and state any access conditions.}

\bibliography{references}

\end{document}
